
\section{OA-03: One round of entanglement distillation}
\source{Week 4 Tutorial, Question 1 (\texttt{Week4TutorialQuestion (1).pdf}).}
\subsection*{Set up a noisy pair}
Alice and Bob share
\[
\rho=p\ket{\Phi^+}\bra{\Phi^+}
+(1-p)\ket{\Psi^+}\bra{\Psi^+},
\qquad \tfrac12<p\le1 .
\]
The input fidelity to \(\ket{\Phi^+}\) is \(F_{\rm in}=p\).
Take two independent copies. Each party performs a CNOT
from their \emph{first} pair to their \emph{second} pair,
then measures the target pair and keeps the control pair
only if the target measurement results agree.

\subsection*{Step 1: classify the four pure branches}
The four possibilities and their probabilities are
\[
\begin{array}{c|c|c}
\text{Control}\otimes\text{target}&\text{Weight}&
\text{Target results agree?}\\\hline
\Phi^+\otimes\Phi^+&p^2&\text{yes}\\
\Phi^+\otimes\Psi^+&p(1-p)&\text{no}\\
\Psi^+\otimes\Phi^+&(1-p)p&\text{no}\\
\Psi^+\otimes\Psi^+&(1-p)^2&\text{yes}
\end{array}
\]
The bilateral CNOT copies the \emph{bit-flip error label}
of the two Bell pairs into the target pair: matching labels
lead to equal target outcomes. For example,
\(\ket{\Phi^+}\ket{\Phi^+}\) remains two \(\Phi^+\) pairs,
whereas \(\ket{\Phi^+}\ket{\Psi^+}\) produces a target
\(\Psi^+\), whose two computational outcomes differ.

\subsection*{Step 2: condition on successful heralding}
The successful event has probability
\[
\boxed{P_{\rm suc}=p^2+(1-p)^2}.
\]
The (unnormalised) retained control state is
\[
\tilde\rho=p^2\ket{\Phi^+}\bra{\Phi^+}
+(1-p)^2\ket{\Psi^+}\bra{\Psi^+}.
\]
Normalise by its trace:
\[
\boxed{\rho'=
\frac{p^2}{p^2+(1-p)^2}\ket{\Phi^+}\bra{\Phi^+}
+\frac{(1-p)^2}{p^2+(1-p)^2}\ket{\Psi^+}\bra{\Psi^+}}.
\]
Therefore
\[
\boxed{F'=\frac{p^2}{p^2+(1-p)^2}>p
\quad\text{for }1/2<p<1}.
\]
For \(p=0.8\), \(P_{\rm suc}=0.68\) and
\(F'=0.64/0.68\simeq0.9412\).

\subsection*{Interpretation}
The fidelity improves \emph{conditional on success}, but
each round consumes two pairs to create at most one.
Repeated rounds improve quality at the expense of yield.
\checkmethod{Postselection requires two distinct quantities:
the trace (success probability) and the normalised
postselected density matrix. Do not confuse them.}
